\section{Localization in categories}
\label{section-localization}

\noindent
The basic idea of this section is given a category $\mathcal{C}$
and a set of arrows $S$ to construct a functor
$F : \mathcal{C} \to S^{-1}\mathcal{C}$
such that all elements of $S$ become invertible in $S^{-1}\mathcal{C}$
and such that $F$ is universal among all functors with this property.
References for this section are \cite[Chapter I, Section 2]{GZ}
and \cite[Chapter II, Section 2]{Verdier}.

\begin{definition}
\label{definition-multiplicative-system}
Let $\mathcal{C}$ be a category. A set of arrows $S$ of $\mathcal{C}$ is
called a {\it left multiplicative system} if it has the following properties:
\begin{enumerate}
\item[LMS1] The identity of every object of $\mathcal{C}$ is in $S$ and
the composition of two composable elements of $S$ is in $S$.
\item[LMS2] Every solid diagram
$$
\xymatrix{
X \ar[d]_t \ar[r]_g & Y \ar@{..>}[d]^s \\
Z \ar@{..>}[r]^f & W
}
$$
with $t \in S$ can be completed to a commutative dotted square with
$s \in S$.
\item[LMS3] For every pair of morphisms $f, g : X \to Y$ and
$t \in S$ with target $X$ such that $f \circ t = g \circ t$
there exists an $s \in S$ with source $Y$ such that
$s \circ f = s \circ g$.
\end{enumerate}
A set of arrows $S$ of $\mathcal{C}$ is
called a {\it right multiplicative system}
if it has the following properties:
\begin{enumerate}
\item[RMS1] The identity of every object of $\mathcal{C}$ is in $S$ and
the composition of two composable elements of $S$ is in $S$.
\item[RMS2] Every solid diagram
$$
\xymatrix{
X \ar@{..>}[d]_t \ar@{..>}[r]_g & Y \ar[d]^s \\
Z \ar[r]^f & W
}
$$
with $s \in S$ can be completed to a commutative dotted square with
$t \in S$.
\item[RMS3] For every pair of morphisms $f, g : X \to Y$ and
$s \in S$ with source $Y$ such that $s \circ f = s \circ g$
there exists a $t \in S$ with target $X$ such that
$f \circ t = g \circ t$.
\end{enumerate}
A set of arrows $S$ of $\mathcal{C}$ is called a {\it multiplicative system}
if it is both a left multiplicative system and a right multiplicative system.
In other words, this means that MS1, MS2, MS3 hold, where
MS1 $=$ LMS1 $+$ RMS1, MS2 $=$ LMS2 $+$ RMS2, and
MS3 $=$ LMS3 $+$ RMS3. (That said, of course LMS1 $=$ RMS1
$=$ MS1.)
\end{definition}

\noindent
These conditions are useful to construct the categories $S^{-1}\mathcal{C}$
as follows.

\medskip\noindent
{\bf Left calculus of fractions.}
Let $\mathcal{C}$ be a category and let $S$ be a left multiplicative
system. We define a new category $S^{-1}\mathcal{C}$ as follows
(we verify this works in the proof of
Lemma \ref{lemma-left-localization}):
\begin{enumerate}
\item We set $\Ob(S^{-1}\mathcal{C}) = \Ob(\mathcal{C})$.
\item Morphisms $X \to Y$ of $S^{-1}\mathcal{C}$ are given by pairs
$(f : X \to Y', s : Y \to Y')$ with $s \in S$ up to equivalence.
(The equivalence is defined below. Think of the equivalence class
of a pair $(f, s)$ as $s^{-1}f : X \to Y$.)
\item Two pairs $(f_1 : X \to Y_1, s_1 : Y \to Y_1)$ and
$(f_2 : X \to Y_2, s_2 : Y \to Y_2)$ are said to be equivalent
if there exist a third pair $(f_3 : X \to Y_3, s_3 : Y \to Y_3)$
and morphisms $u : Y_1 \to Y_3$ and $v : Y_2 \to Y_3$ of $\mathcal{C}$
fitting into the commutative diagram
$$
\xymatrix{
 & Y_1 \ar[d]^u & \\
X \ar[ru]^{f_1} \ar[r]^{f_3} \ar[rd]_{f_2} &
Y_3 &
Y \ar[lu]_{s_1} \ar[l]_{s_3} \ar[ld]^{s_2} \\
& Y_2 \ar[u]_v &
}
$$
\item The composition of the equivalence classes of the pairs
$(f : X \to Y', s : Y \to Y')$ and $(g : Y \to Z', t : Z \to Z')$
is defined as the equivalence class of a pair
$(h \circ f : X \to Z'', u \circ t : Z \to Z'')$
where $h$ and $u \in S$ are chosen to fit into a commutative diagram
$$
\xymatrix{
Y \ar[d]_s \ar[r]_g & Z' \ar[d]^u \\
Y' \ar[r]^h & Z''
}
$$
which exists by assumption.
\item The identity morphism $X \to X$ in $S^{-1} \mathcal{C}$ is the
equivalence class of the pair $(\text{id} : X \to X,
\text{id} : X \to X)$.
\end{enumerate}

\begin{lemma}
\label{lemma-left-localization}
Let $\mathcal{C}$ be a category and let $S$ be a left multiplicative
system.
\begin{enumerate}
\item The relation on pairs defined above is an equivalence relation.
\item The composition rule given above is well defined on equivalence
classes.
\item Composition is associative (and the identity morphisms satisfy
the identity axioms), and hence $S^{-1}\mathcal{C}$ is a category.
\end{enumerate}
\end{lemma}

\begin{proof}
Proof of (1). Let us say two pairs $p_1 = (f_1 : X \to Y_1, s_1 : Y \to Y_1)$
and $p_2 = (f_2 : X \to Y_2, s_2 : Y \to Y_2)$ are elementary equivalent
if there exists a morphism $a : Y_1 \to Y_2$ of $\mathcal{C}$ such that
$a \circ f_1 = f_2$ and $a \circ s_1 = s_2$. Diagram:
$$
\xymatrix{
X \ar@{=}[d] \ar[r]_{f_1} & Y_1 \ar[d]^a & Y \ar[l]^{s_1} \ar@{=}[d] \\
X \ar[r]^{f_2} & Y_2 & Y \ar[l]_{s_2}
}
$$
Let us denote this property by saying $p_1Ep_2$.
Note that $pEp$ and $aEb, bEc \Rightarrow aEc$.
(Despite its name, $E$ is not an equivalence
relation.)
Part (1) claims that the relation
$p \sim p' \Leftrightarrow \exists q: pEq \wedge p'Eq$
(where $q$ is supposed to be a pair satisfying the
same conditions as $p$ and $p'$)
is an equivalence relation. A simple formal argument, using the properties
of $E$ above, shows that it suffices to prove
$p_3Ep_1, p_3Ep_2 \Rightarrow p_1 \sim p_2$.
Thus suppose that we are given a commutative diagram
$$
\xymatrix{
 & Y_1 & \\
X \ar[ru]^{f_1} \ar[r]^{f_3} \ar[rd]_{f_2} &
Y_3 \ar[u]_{a_{31}} \ar[d]^{a_{32}} &
Y \ar[lu]_{s_1} \ar[l]_{s_3} \ar[ld]^{s_2} \\
& Y_2 &
}
$$
with $s_i \in S$.
First we apply LMS2 to get a commutative diagram
$$
\xymatrix{
Y \ar[d]_{s_1} \ar[r]_{s_2} & Y_2 \ar@{..>}[d]^{a_{24}} \\
Y_1 \ar@{..>}[r]^{a_{14}} & Y_4
}
$$
with $a_{24} \in S$. Then, we have
$$
a_{14} \circ a_{31} \circ s_3 =
a_{14} \circ s_1 =
a_{24} \circ s_2 =
a_{24} \circ a_{32} \circ s_3.
$$
Hence, by LMS3, there exists a
morphism $s_{44} : Y_4 \to Y'_4$ such that $s_{44} \in S$ and
$s_{44} \circ a_{14} \circ a_{31}
= s_{44} \circ a_{24} \circ a_{32}$.
Hence, after replacing $Y_4$, $a_{14}$ and $a_{24}$ by $Y'_4$,
$s_{44} \circ a_{14}$ and $s_{44} \circ a_{24}$, we may assume
that $a_{14} \circ a_{31} = a_{24} \circ a_{32}$ (and
we still have $a_{24} \in S$ and
$a_{14} \circ s_1 = a_{24} \circ s_2$). Set
$$
f_4 =
a_{14} \circ f_1 =
a_{14} \circ a_{31} \circ f_3 =
a_{24} \circ a_{32} \circ f_3 =
a_{24} \circ f_2
$$
and
$s_4 = a_{14} \circ s_1 = a_{24} \circ s_2$. Then, the diagram
$$
\xymatrix{
X \ar@{=}[d] \ar[r]_{f_1} & Y_1 \ar[d]^{a_{14}} & Y \ar[l]^{s_1} \ar@{=}[d] \\
X \ar[r]^{f_4} & Y_4 & Y \ar[l]_{s_4}
}
$$
commutes, and we have $s_4 \in S$ (by LMS1). Thus, $p_1 E p_4$,
where $p_4 = (f_4, s_4)$. Similarly, $p_2 E p_4$. Combining these,
we find $p_1 \sim p_2$.

\medskip\noindent
Proof of (2). Let $p = (f : X \to Y', s : Y \to Y')$ and
$q = (g : Y \to Z', t : Z \to Z')$ be pairs as in the definition of composition
above. To compose we choose a diagram
$$
\xymatrix{
Y \ar[d]_s \ar[r]_g & Z' \ar[d]^{u_2} \\
Y' \ar[r]^{h_2} & Z_2
}
$$
with $u_2 \in S$. We first show that the equivalence class of the pair
$r_2 = (h_2 \circ f : X \to Z_2, u_2 \circ t : Z \to Z_2)$
is independent of the choice of $(Z_2, h_2, u_2)$. Namely, suppose
that $(Z_3, h_3, u_3)$ is another choice with corresponding composition
$r_3 = (h_3 \circ f : X \to Z_3, u_3 \circ t : Z \to Z_3)$.
Then by LMS2 we can choose a diagram
$$
\xymatrix{
Z' \ar[d]_{u_2} \ar[r]_{u_3} & Z_3 \ar[d]^{u_{34}} \\
Z_2 \ar[r]^{h_{24}} & Z_4
}
$$
with $u_{34} \in S$. We have $h_2 \circ s = u_2 \circ g$
and similarly $h_3 \circ s = u_3 \circ g$. Now,
$$
u_{34} \circ h_3 \circ s
= u_{34} \circ u_3 \circ g
= h_{24} \circ u_2 \circ g
= h_{24} \circ h_2 \circ s.
$$
Hence, LMS3 shows that there
exist a $Z'_4$ and an $s_{44} : Z_4 \to Z'_4$ such that
$s_{44} \circ u_{34} \circ h_3 = s_{44} \circ h_{24} \circ h_2$.
Replacing $Z_4$, $h_{24}$ and $u_{34}$ by $Z'_4$,
$s_{44} \circ h_{24}$ and $s_{44} \circ u_{34}$, we may
assume that $u_{34} \circ h_3 = h_{24} \circ h_2$.
Meanwhile, the relations $u_{34} \circ u_3 = h_{24} \circ u_2$
and $u_{34} \in S$ continue to hold. We can now set
$h_4 = u_{34} \circ h_3 = h_{24} \circ h_2$ and
$u_4 = u_{34} \circ u_3 = h_{24} \circ u_2$. Then, we have a
commutative diagram
$$
\xymatrix{
X \ar@{=}[d] \ar[r]_{h_2\circ f} &
Z_2 \ar[d]^{h_{24}} &
Z \ar[l]^{u_2 \circ t} \ar@{=}[d] \\
X \ar@{=}[d] \ar[r]^{h_4\circ f} &
Z_4 &
Z \ar@{=}[d] \ar[l]_{u_4 \circ t} \\
X \ar[r]^{h_3 \circ f} &
Z_3 \ar[u]^{u_{34}} &
Z \ar[l]_{u_3 \circ t}
}
$$
Hence we obtain a pair
$r_4 =
(h_4 \circ f : X \to Z_4, u_4 \circ t : Z \to Z_4)$
and the above diagram shows that we have $r_2Er_4$ and
$r_3Er_4$, whence $r_2 \sim r_3$, as desired. Thus it now
makes sense to define $p \circ q$ as the equivalence class of
all possible pairs $r$ obtained as above.

\medskip\noindent
To finish the proof of (2) we have to show that given pairs
$p_1, p_2, q$ such that $p_1Ep_2$ then $p_1 \circ q = p_2 \circ q$ and
$q \circ p_1 = q \circ p_2$ whenever the compositions make sense.
To do this, write $p_1 = (f_1 : X \to Y_1, s_1 : Y \to Y_1)$ and
$p_2 = (f_2 : X \to Y_2, s_2 : Y \to Y_2)$ and let
$a : Y_1 \to Y_2$ be a morphism of $\mathcal{C}$ such that
$f_2 = a \circ f_1$ and $s_2 = a \circ s_1$.
First assume that $q = (g : Y \to Z', t : Z \to Z')$.
In this case choose a commutative diagram as the one on the left
$$
\vcenter{
\xymatrix{
Y \ar[d]_{s_2} \ar[r]^g & Z' \ar[d]^u \\
Y_2 \ar[r]^h & Z''
}
}
\quad
\Rightarrow
\quad
\vcenter{
\xymatrix{
Y \ar[d]_{s_1} \ar[r]^g & Z' \ar[d]^u \\
Y_1 \ar[r]^{h \circ a} & Z''
}
}
$$
(with $u \in S$),
which implies the diagram on the right is commutative as well.
Using these diagrams we see that both compositions $q \circ p_1$
and $q \circ p_2$ are the equivalence class of
$(h \circ a \circ f_1 : X \to Z'', u \circ t : Z \to Z'')$.
Thus $q \circ p_1 = q \circ p_2$.
The proof of the other case, in which we have to show
$p_1 \circ q = p_2 \circ q$, is omitted. (It is similar to the
case we did.)

\medskip\noindent
Proof of (3). We have to prove associativity of composition.
Consider a solid diagram
$$
\xymatrix{
& & & Z \ar[d] \\
& & Y \ar[d] \ar[r] & Z' \ar@{..>}[d] \\
& X \ar[d] \ar[r] & Y' \ar@{..>}[d] \ar@{..>}[r] & Z'' \ar@{..>}[d] \\
W \ar[r] & X' \ar@{..>}[r] & Y'' \ar@{..>}[r] & Z'''
}
$$
(whose vertical arrows belong to $S$)
which gives rise to three composable pairs.
Using LMS2 we can choose the dotted arrows making the squares commutative
and such that the vertical arrows are in $S$.
Then it is clear that the composition of the three pairs
is the equivalence class of the pair
$(W \to Z''', Z \to Z''')$ gotten by composing the
horizontal arrows on the bottom row and the vertical arrows
on the right column.

\medskip\noindent
We leave it to the reader to check the identity axioms.
\end{proof}

\begin{remark}
\label{remark-motivation-localization}
The motivation for the construction of $S^{-1} \mathcal{C}$ is to
``force'' the morphisms in $S$ to be invertible by artificially
creating inverses to them (at the cost of some existing morphisms
possibly becoming identified with each other). This is similar to
the localization of a commutative ring at a multiplicative subset,
and more generally to the localization of a noncommutative ring
at a right denominator set (see \cite[Section 10A]{Lam}). This is
more than just a similarity: The construction of
$S^{-1} \mathcal{C}$ (or, more precisely, its version for
additive categories $\mathcal{C}$) actually generalizes the
latter type of localization. Namely, a noncommutative ring can be
viewed as a pre-additive category with a single object (the morphisms
being the elements of the ring); a multiplicative subset of this
ring then becomes a set $S$ of morphisms satisfying LMS1 (aka
RMS1). Then, the conditions RMS2 and RMS3 for this category and
this subset $S$ translate into the two conditions
(``right permutable'' and ``right reversible'') of a right
denominator set (and similarly for LMS and left denominator sets),
and $S^{-1} \mathcal{C}$ (with a properly defined additive
structure) is the one-object category corresponding to the
localization of the ring.
\end{remark}

\begin{definition}
\label{definition-left-localization-as-fraction}
Let $\mathcal{C}$ be a category and let $S$ be a left multiplicative
system of morphisms of $\mathcal{C}$. Given any morphism
$f : X \to Y'$ in $\mathcal{C}$ and any morphism $s : Y \to Y'$ in
$S$, we denote by {\it $s^{-1} f$} the equivalence class of the pair
$(f : X \to Y', s : Y \to Y')$. This is a morphism from $X$ to $Y$
in $S^{-1} \mathcal{C}$.
\end{definition}

\noindent
This notation is suggestive, and the things it suggests are true:
Given any morphism $f : X \to Y'$ in $\mathcal{C}$ and any two
morphisms $s : Y \to Y'$ and $t : Y' \to Y''$ in $S$, we have
$\left(t \circ s\right)^{-1} \left(t \circ f\right) = s^{-1} f$.
Also, for any
$f : X \to Y'$ and $g : Y' \to Z'$ in $\mathcal{C}$ and all
$s : Z \to Z'$ in $S$, we have
$s^{-1} \left(g \circ f\right) = \left(s^{-1} g\right) \circ
\left(\text{id}_{Y'}^{-1} f\right)$.
Finally, for any $f : X \to Y'$ in $\mathcal{C}$, all
$s : Y \to Y'$ in $S$, and $t : Z \to Y$ in $S$, we have
$\left(s \circ t\right)^{-1} f
= \left(t^{-1} \text{id}_Y\right)
\circ \left(s^{-1} f\right)$.
This is all clear from the definition.
We can ``write any finite collection of morphisms with the same target
as fractions with common denominator''.

\begin{lemma}
\label{lemma-morphisms-left-localization}
Let $\mathcal{C}$ be a category and let $S$ be a left multiplicative
system of morphisms of $\mathcal{C}$. Given any finite collection
$g_i : X_i \to Y$ of morphisms of $S^{-1}\mathcal{C}$
(indexed by $i$),
we can find an element $s : Y \to Y'$ of $S$ and
a family of morphisms $f_i : X_i \to Y'$ of $\mathcal{C}$ such that
each $g_i$ is the equivalence class of the pair
$(f_i : X_i \to Y', s : Y \to Y')$.
\end{lemma}

\begin{proof}
For each $i$ choose a representative $(X_i \to Y_i, s_i : Y \to Y_i)$
of $g_i$.
The lemma follows if we can find a morphism $s : Y \to Y'$ in $S$ such that
for each $i$ there is a morphism $a_i : Y_i \to Y'$ with
$a_i \circ s_i = s$. If we have two indices $i = 1, 2$, then we can
do this by completing the square
$$
\xymatrix{
Y \ar[d]_{s_1} \ar[r]_{s_2} & Y_2 \ar[d]^{t_2} \\
Y_1 \ar[r]^{a_1} & Y'
}
$$
with $t_2 \in S$ as is possible by
Definition \ref{definition-multiplicative-system}.
Then $s = t_2 \circ s_2 \in S$ works.
If we have $n > 2$ morphisms, then we use the above trick to reduce
to the case of $n - 1$ morphisms, and we win by induction.
\end{proof}

\noindent
There is an easy characterization of equality of morphisms if they
have the same denominator.

\begin{lemma}
\label{lemma-equality-morphisms-left-localization}
Let $\mathcal{C}$ be a category and let $S$ be a left multiplicative
system of morphisms of $\mathcal{C}$. Let $A, B : X \to Y$ be morphisms
of $S^{-1}\mathcal{C}$ which are the equivalence classes of
$(f : X \to Y', s : Y \to Y')$ and $(g : X \to Y', s : Y \to Y')$.
The following are equivalent
\begin{enumerate}
\item $A = B$
\item there exists a morphism $t : Y' \to Y''$
in $S$ with $t \circ f = t \circ g$, and
\item there exists a morphism $a : Y' \to Y''$
such that $a \circ f = a \circ g$ and $a \circ s \in S$.
\end{enumerate}
\end{lemma}

\begin{proof}
We are going to use that $S^{-1}\mathcal{C}$ is a category
(Lemma \ref{lemma-left-localization}) and we will
use the notation of Definition \ref{definition-left-localization-as-fraction}
as well as the discussion following that definition to identify some
morphisms in $S^{-1}\mathcal{C}$. Thus we write $A = s^{-1}f$ and $B = s^{-1}g$.

\medskip\noindent
If $A = B$ then
$(\text{id}_{Y'}^{-1}s) \circ A = (\text{id}_{Y'}^{-1}s) \circ B$.
We have $(\text{id}_{Y'}^{-1}s) \circ A = \text{id}_{Y'}^{-1}f$
and $(\text{id}_{Y'}^{-1}s) \circ B = \text{id}_{Y'}^{-1}g$.
The equality of $\text{id}_{Y'}^{-1}f$ and $\text{id}_{Y'}^{-1}g$
means by definition that there exists a commutative diagram
$$
\xymatrix{
 & Y' \ar[d]^u & \\
X \ar[ru]^f \ar[r]^h \ar[rd]_g &
Z &
Y' \ar[lu]_{\text{id}_{Y'}} \ar[l]_t \ar[ld]^{\text{id}_{Y'}} \\
& Y' \ar[u]_v &
}
$$
with $t \in S$. In particular $u = v = t \in S$ and $t \circ f = t\circ g$.
Thus (1) implies (2).

\medskip\noindent
The implication (2) $\Rightarrow$ (3) is immediate. Assume $a$ is as in (3).
Denote $s' = a \circ s \in S$.
Then $\text{id}_{Y''}^{-1}s'$ is an isomorphism in the category
$S^{-1}\mathcal{C}$ (with inverse $(s')^{-1}\text{id}_{Y''}$).
Thus to check $A = B$ it suffices to check that
$\text{id}_{Y''}^{-1}s' \circ A = \text{id}_{Y''}^{-1}s' \circ B$.
We compute using the rules discussed in the text following
Definition \ref{definition-left-localization-as-fraction} that
$\text{id}_{Y''}^{-1}s' \circ A =
\text{id}_{Y''}^{-1}(a \circ s) \circ s^{-1}f =
\text{id}_{Y''}^{-1}(a \circ f) =
\text{id}_{Y''}^{-1}(a \circ g) =
\text{id}_{Y''}^{-1}(a \circ s) \circ s^{-1}g =
\text{id}_{Y''}^{-1}s' \circ B$ and we see that (1) is true.
\end{proof}

\begin{remark}
\label{remark-left-localization-morphisms-colimit}
Let $\mathcal{C}$ be a category. Let $S$ be a left multiplicative system.
Given an object $Y$ of $\mathcal{C}$ we denote $Y/S$ the category whose
objects are $s : Y \to Y'$ with $s \in S$ and whose morphisms are
commutative diagrams
$$
\xymatrix{
& Y \ar[ld]_s \ar[rd]^t & \\
Y' \ar[rr]^a & & Y''
}
$$
where $a : Y' \to Y''$ is arbitrary. We claim that the category
$Y/S$ is filtered (see
Definition \ref{definition-directed}).
Namely, LMS1 implies that $\text{id}_Y : Y \to Y$
is in $Y/S$; hence $Y/S$ is nonempty. LMS2 implies that given
$s_1 : Y \to Y_1$ and $s_2 : Y \to Y_2$ we can find a diagram
$$
\xymatrix{
Y \ar[d]_{s_1} \ar[r]_{s_2} & Y_2 \ar[d]^t \\
Y_1 \ar[r]^a & Y_3
}
$$
with $t \in S$. Hence $s_1 : Y \to Y_1$ and $s_2 : Y \to Y_2$
both have maps to $t \circ s_2 : Y \to Y_3$ in $Y/S$. Finally, given
two morphisms $a, b$ from $s_1 : Y \to Y_1$ to $s_2 : Y \to Y_2$
in $Y/S$ we see that $a \circ s_1 = b \circ s_1$; hence by LMS3
there exists a $t : Y_2 \to Y_3$ in $S$ such that
$t \circ a = t \circ b$.
Now the combined results of
Lemmas \ref{lemma-morphisms-left-localization} and
\ref{lemma-equality-morphisms-left-localization}
tell us that
\begin{equation}
\label{equation-left-localization-morphisms-colimit}
\Mor_{S^{-1}\mathcal{C}}(X, Y) =
\colim_{(s : Y \to Y') \in Y/S} \Mor_\mathcal{C}(X, Y')
\end{equation}
This formula expressing morphism sets in $S^{-1}\mathcal{C}$ as a filtered
colimit of morphism sets in $\mathcal{C}$ is occasionally useful.
\end{remark}

\begin{lemma}
\label{lemma-properties-left-localization}
Let $\mathcal{C}$ be a category and let $S$ be a left multiplicative
system of morphisms of $\mathcal{C}$.
\begin{enumerate}
\item The rules $X \mapsto X$ and
$(f : X \to Y) \mapsto (f : X \to Y, \text{id}_Y : Y \to Y)$
define a functor $Q : \mathcal{C} \to S^{-1}\mathcal{C}$.
\item For any $s \in S$ the morphism $Q(s)$ is an isomorphism in
$S^{-1}\mathcal{C}$.
\item If $G : \mathcal{C} \to \mathcal{D}$ is any functor such that
$G(s)$ is invertible for every $s \in S$, then there exists a
unique functor $H : S^{-1}\mathcal{C} \to \mathcal{D}$
such that $H \circ Q = G$.
\end{enumerate}
\end{lemma}

\begin{proof}
Parts (1) and (2) are clear. (In (2), the inverse of $Q(s)$ is
the equivalence class of the pair $(\text{id}_Y, s)$.)
To see (3) just set $H(X) = G(X)$
and set $H((f : X \to Y', s : Y \to Y')) = G(s)^{-1} \circ G(f)$.
Details omitted.
\end{proof}

\begin{lemma}
\label{lemma-left-localization-limits}
Let $\mathcal{C}$ be a category and let $S$ be a left multiplicative
system of morphisms of $\mathcal{C}$. The localization functor
$Q : \mathcal{C} \to S^{-1}\mathcal{C}$ commutes with finite colimits.
\end{lemma}

\begin{proof}
Let $\mathcal{I}$ be a finite category and let
$\mathcal{I} \to \mathcal{C}$, $i \mapsto X_i$
be a functor whose colimit exists. Then using
(\ref{equation-left-localization-morphisms-colimit}),
the fact that $Y/S$ is filtered, and
Lemma \ref{lemma-directed-commutes} we have
\begin{align*}
\Mor_{S^{-1}\mathcal{C}}(Q(\colim X_i), Q(Y))
& =
\colim_{(s : Y \to Y') \in Y/S} \Mor_\mathcal{C}(\colim X_i, Y') \\
& =
\colim_{(s : Y \to Y') \in Y/S} \lim_i \Mor_\mathcal{C}(X_i, Y') \\
& =
\lim_i \colim_{(s : Y \to Y') \in Y/S} \Mor_\mathcal{C}(X_i, Y') \\
& =
\lim_i \Mor_{S^{-1}\mathcal{C}}(Q(X_i), Q(Y))
\end{align*}
and this isomorphism commutes with the projections
from both sides to the set
$\Mor_{S^{-1}\mathcal{C}}(Q(X_j), Q(Y))$ for each
$j \in \Ob(\mathcal{I})$. Thus, $Q(\colim X_i)$ satisfies
the universal property for the colimit of the functor
$i \mapsto Q(X_i)$; hence, it is this colimit, as desired.
\end{proof}

\begin{lemma}
\label{lemma-left-localization-diagram}
Let $\mathcal{C}$ be a category. Let $S$ be a left multiplicative
system. If $f : X \to Y$, $f' : X' \to Y'$ are two morphisms of
$\mathcal{C}$ and if
$$
\xymatrix{
Q(X) \ar[d]_{Q(f)} \ar[r]_a & Q(X') \ar[d]^{Q(f')} \\
Q(Y) \ar[r]^b & Q(Y')
}
$$
is a commutative diagram in $S^{-1}\mathcal{C}$, then there exist
a morphism $f'' : X'' \to Y''$ in $\mathcal{C}$ and a commutative
diagram
$$
\xymatrix{
X \ar[d]_f \ar[r]_g & X'' \ar[d]^{f''} & X' \ar[d]^{f'} \ar[l]^s \\
Y \ar[r]^h & Y'' & Y' \ar[l]_t
}
$$
in $\mathcal{C}$ with $s, t \in S$ and $a = s^{-1}g$, $b = t^{-1}h$.
\end{lemma}

\begin{proof}
We choose maps and objects in the following way:
First write $a = s^{-1}g$ for some $s : X' \to X''$ in $S$ and
$g : X \to X''$. By LMS2 we can find $t : Y' \to Y''$ in $S$ and
$f'' : X'' \to Y''$ such that
$$
\xymatrix{
X \ar[d]_f \ar[r]_g & X'' \ar[d]^{f''} & X' \ar[d]^{f'} \ar[l]^s \\
Y & Y'' & Y' \ar[l]_t
}
$$
commutes. Now in this diagram we are going to repeatedly change our
choice of
$$
X'' \xrightarrow{f''} Y'' \xleftarrow{t} Y'
$$
by postcomposing both $t$ and $f''$ by a morphism $d : Y'' \to Y'''$
with the property that $d \circ t \in S$. According to
Remark \ref{remark-left-localization-morphisms-colimit}
we may after such a replacement assume that there exists a morphism
$h : Y \to Y''$ such that $b = t^{-1}h$ holds\footnote{Here is a
more down-to-earth way to see this:
Write $b = q^{-1}i$ for some $q : Y' \to Z$ in $S$ and some
$i : Y \to Z$. By LMS2 we can find $r : Y'' \to Y'''$ in $S$ and
$j : Z \to Y'''$ such that $j \circ q = r \circ t$. Now, set
$d = r$ and $h = j \circ i$.}. At this point we have everything
as in the lemma except that we don't know that the left square of the
diagram commutes.
But the definition of composition in $S^{-1} \mathcal{C}$ shows that
$b \circ Q\left(f\right)$ is the equivalence class of the pair
$(h \circ f : X \to Y'', t : Y' \to Y'')$ (since $b$ is the
equivalence class of the pair $(h : Y \to Y'', t : Y' \to Y'')$,
while $Q\left(f\right)$ is the equivalence class of the pair
$(f : X \to Y, \text{id} : Y \to Y)$), while
$Q\left(f'\right) \circ a$ is the equivalence class of the pair
$(f'' \circ g : X \to Y'', t : Y' \to Y'')$ (since $a$ is the
equivalence class of the pair $(g : X \to X'', s : X' \to X'')$,
while $Q\left(f'\right)$ is the equivalence class of the pair
$(f' : X' \to Y', \text{id} : Y' \to Y')$).
Since we know that
$b \circ Q\left(f\right) = Q\left(f'\right) \circ a$, we thus
conclude that the equivalence classes of the pairs
$(h \circ f : X \to Y'', t : Y' \to Y'')$ and
$(f'' \circ g : X \to Y'', t : Y' \to Y'')$ are equal.
Hence using
Lemma \ref{lemma-equality-morphisms-left-localization}
we can find a morphism $d : Y'' \to Y'''$ such that
$d \circ t \in S$ and $d \circ h \circ f = d \circ f'' \circ g$.
Hence we make one more replacement of the kind described
above and we win.
\end{proof}

\noindent
{\bf Right calculus of fractions.}
Let $\mathcal{C}$ be a category and let $S$ be a right multiplicative
system. We define a new category $S^{-1}\mathcal{C}$ as follows
(we verify this works in the proof of
Lemma \ref{lemma-right-localization}):
\begin{enumerate}
\item We set $\Ob(S^{-1}\mathcal{C}) = \Ob(\mathcal{C})$.
\item Morphisms $X \to Y$ of $S^{-1}\mathcal{C}$ are given by pairs
$(f : X' \to Y, s : X' \to X)$ with $s \in S$ up to equivalence.
(The equivalence is defined below. Think of the equivalence class
of a pair $(f, s)$ as $fs^{-1} : X \to Y$.)
\item Two pairs $(f_1 : X_1 \to Y, s_1 : X_1 \to X)$ and
$(f_2 : X_2 \to Y, s_2 : X_2 \to X)$ are said to be equivalent
if there exist a third pair $(f_3 : X_3 \to Y, s_3 : X_3 \to X)$
and morphisms $u : X_3 \to X_1$ and $v : X_3 \to X_2$ of $\mathcal{C}$
fitting into the commutative diagram
$$
\xymatrix{
 & X_1 \ar[ld]_{s_1} \ar[rd]^{f_1} & \\
X &
X_3 \ar[l]_{s_3} \ar[u]_u \ar[d]^v \ar[r]^{f_3} &
Y \\
& X_2 \ar[lu]^{s_2} \ar[ru]_{f_2} &
}
$$
\item The composition of the equivalence classes of the pairs
$(f : X' \to Y, s : X' \to X)$ and $(g : Y' \to Z, t : Y' \to Y)$
is defined as the equivalence class of a pair
$(g \circ h : X'' \to Z, s \circ u : X'' \to X)$
where $h$ and $u \in S$ are chosen to fit into a commutative diagram
$$
\xymatrix{
X'' \ar[d]_u \ar[r]^h & Y' \ar[d]^t \\
X' \ar[r]^f & Y
}
$$
which exists by assumption.
\item The identity morphism $X \to X$ in $S^{-1} \mathcal{C}$ is the
equivalence class of the pair $(\text{id} : X \to X,
\text{id} : X \to X)$.
\end{enumerate}

\begin{lemma}
\label{lemma-right-localization}
Let $\mathcal{C}$ be a category and let $S$ be a right multiplicative
system.
\begin{enumerate}
\item The relation on pairs defined above is an equivalence relation.
\item The composition rule given above is well defined on equivalence
classes.
\item Composition is associative (and the identity morphisms satisfy
the identity axioms), and hence $S^{-1}\mathcal{C}$ is a category.
\end{enumerate}
\end{lemma}

\begin{proof}
This lemma is dual to
Lemma \ref{lemma-left-localization}.
It follows formally from that lemma by replacing
$\mathcal{C}$ by its opposite category in which
$S$ is a left multiplicative system.
\end{proof}

\begin{definition}
\label{definition-right-localization-as-fraction}
Let $\mathcal{C}$ be a category and let $S$ be a right multiplicative
system of morphisms of $\mathcal{C}$. Given any morphism
$f : X' \to Y$ in $\mathcal{C}$ and any morphism $s : X' \to X$ in
$S$, we denote by {\it $f s^{-1}$} the equivalence class of the pair
$(f : X' \to Y, s : X' \to X)$. This is a morphism from $X$ to $Y$
in $S^{-1} \mathcal{C}$.
\end{definition}

\noindent
Identities similar (actually, dual) to the ones in Definition
\ref{definition-left-localization-as-fraction} hold.
We can ``write any finite collection of morphisms with the same source
as fractions with common denominator''.

\begin{lemma}
\label{lemma-morphisms-right-localization}
Let $\mathcal{C}$ be a category and let $S$ be a right multiplicative
system of morphisms of $\mathcal{C}$. Given any finite collection
$g_i : X \to Y_i$ of morphisms of $S^{-1}\mathcal{C}$
(indexed by $i$),
we can find an element $s : X' \to X$ of $S$ and a family
of morphisms $f_i : X' \to Y_i$ of $\mathcal{C}$ such that
$g_i$ is the equivalence class of the pair
$(f_i : X' \to Y_i, s : X' \to X)$.
\end{lemma}

\begin{proof}
This lemma is the dual of
Lemma \ref{lemma-morphisms-left-localization}
and follows formally from that lemma by replacing all
categories in sight by their opposites.
\end{proof}

\noindent
There is an easy characterization of equality of morphisms if they
have the same denominator.

\begin{lemma}
\label{lemma-equality-morphisms-right-localization}
Let $\mathcal{C}$ be a category and let $S$ be a right multiplicative
system of morphisms of $\mathcal{C}$. Let $A, B : X \to Y$ be
morphisms of $S^{-1}\mathcal{C}$ which are the equivalence
classes of $(f : X' \to Y, s : X' \to X)$ and
$(g : X' \to Y, s : X' \to X)$. The following are equivalent
\begin{enumerate}
\item $A = B$,
\item there exists a morphism $t : X'' \to X'$ in $S$ with
$f \circ t = g \circ t$, and
\item there exists a morphism $a : X'' \to X'$ with
$f \circ a = g \circ a$ and $s \circ a \in S$.
\end{enumerate}
\end{lemma}

\begin{proof}
This is dual to
Lemma \ref{lemma-equality-morphisms-left-localization}.
\end{proof}

\begin{remark}
\label{remark-right-localization-morphisms-colimit}
Let $\mathcal{C}$ be a category. Let $S$ be a right multiplicative system.
Given an object $X$ of $\mathcal{C}$ we denote $S/X$ the category whose
objects are $s : X' \to X$ with $s \in S$ and whose morphisms are
commutative diagrams
$$
\xymatrix{
X' \ar[rd]_s \ar[rr]_a & & X'' \ar[ld]^t \\
& X
}
$$
where $a : X' \to X''$ is arbitrary. The category
$S/X$ is cofiltered (see
Definition \ref{definition-codirected}).
(This is dual to the corresponding statement in
Remark \ref{remark-left-localization-morphisms-colimit}.)
Now the combined results of
Lemmas \ref{lemma-morphisms-right-localization} and
\ref{lemma-equality-morphisms-right-localization}
tell us that
\begin{equation}
\label{equation-right-localization-morphisms-colimit}
\Mor_{S^{-1}\mathcal{C}}(X, Y) =
\colim_{(s : X' \to X) \in (S/X)^{opp}} \Mor_\mathcal{C}(X', Y)
\end{equation}
This formula expressing morphisms in $S^{-1}\mathcal{C}$ as a filtered
colimit of morphisms in $\mathcal{C}$ is occasionally useful.
\end{remark}

\begin{lemma}
\label{lemma-properties-right-localization}
Let $\mathcal{C}$ be a category and let $S$ be a right multiplicative
system of morphisms of $\mathcal{C}$.
\begin{enumerate}
\item The rules $X \mapsto X$ and
$(f : X \to Y) \mapsto (f : X \to Y, \text{id}_X : X \to X)$
define a functor $Q : \mathcal{C} \to S^{-1}\mathcal{C}$.
\item For any $s \in S$ the morphism $Q(s)$ is an isomorphism in
$S^{-1}\mathcal{C}$.
\item If $G : \mathcal{C} \to \mathcal{D}$ is any functor such that
$G(s)$ is invertible for every $s \in S$, then there exists a
unique functor $H : S^{-1}\mathcal{C} \to \mathcal{D}$
such that $H \circ Q = G$.
\end{enumerate}
\end{lemma}

\begin{proof}
This lemma is the dual of
Lemma \ref{lemma-properties-left-localization}
and follows formally from that lemma by replacing all
categories in sight by their opposites.
\end{proof}

\begin{lemma}
\label{lemma-right-localization-limits}
Let $\mathcal{C}$ be a category and let $S$ be a right multiplicative
system of morphisms of $\mathcal{C}$. The localization functor
$Q : \mathcal{C} \to S^{-1}\mathcal{C}$ commutes with finite limits.
\end{lemma}

\begin{proof}
This is dual to Lemma \ref{lemma-left-localization-limits}.
\end{proof}

\begin{lemma}
\label{lemma-right-localization-diagram}
Let $\mathcal{C}$ be a category. Let $S$ be a right multiplicative
system. If $f : X \to Y$, $f' : X' \to Y'$ are two morphisms of
$\mathcal{C}$ and if
$$
\xymatrix{
Q(X) \ar[d]_{Q(f)} \ar[r]_a & Q(X') \ar[d]^{Q(f')} \\
Q(Y) \ar[r]^b & Q(Y')
}
$$
is a commutative diagram in $S^{-1}\mathcal{C}$, then there exist
a morphism $f'' : X'' \to Y''$ in $\mathcal{C}$ and a commutative
diagram
$$
\xymatrix{
X \ar[d]_f & X'' \ar[l]^s \ar[d]^{f''} \ar[r]_g & X' \ar[d]^{f'} \\
Y & Y'' \ar[l]_t \ar[r]^h & Y'
}
$$
in $\mathcal{C}$ with $s, t \in S$ and $a = gs^{-1}$, $b = ht^{-1}$.
\end{lemma}

\begin{proof}
This lemma is dual to
Lemma \ref{lemma-left-localization-diagram}.
\end{proof}

\noindent
{\bf Multiplicative systems and two sided calculus of fractions.}
If $S$ is a multiplicative system then left and right calculus of
fractions give canonically isomorphic categories.

\begin{lemma}
\label{lemma-multiplicative-system}
Let $\mathcal{C}$ be a category and let $S$ be a multiplicative system.
The category of left fractions and the category of right fractions
$S^{-1}\mathcal{C}$ are canonically isomorphic.
\end{lemma}

\begin{proof}
Denote $\mathcal{C}_{left}$, $\mathcal{C}_{right}$ the two categories
of fractions. By the universal properties of
Lemmas \ref{lemma-properties-left-localization} and
\ref{lemma-properties-right-localization}
we obtain functors $\mathcal{C}_{left} \to \mathcal{C}_{right}$
and $\mathcal{C}_{right} \to \mathcal{C}_{left}$.
By the uniqueness statement in the universal properties, these
functors are each other's inverse.
\end{proof}

\begin{definition}
\label{definition-saturated-multiplicative-system}
Let $\mathcal{C}$ be a category and let $S$ be a multiplicative system.
We say $S$ is {\it saturated} if, in addition to MS1, MS2, MS3, we
also have
\begin{enumerate}
\item[MS4] Given three composable morphisms $f, g, h$, if
$fg, gh \in S$, then $g \in S$.
\end{enumerate}
\end{definition}

\noindent
Note that a saturated multiplicative system contains all isomorphisms.
Moreover, if $f, g, h$ are composable morphisms in a category and
$fg, gh$ are isomorphisms, then $g$ is an isomorphism (because then $g$
has both a left and a right inverse, hence is invertible).

\begin{lemma}
\label{lemma-what-gets-inverted}
Let $\mathcal{C}$ be a category and let $S$ be a multiplicative system.
Denote $Q : \mathcal{C} \to S^{-1}\mathcal{C}$ the localization functor.
The set
$$
\hat S = \{f \in \text{Arrows}(\mathcal{C}) \mid
Q(f) \text{ is an isomorphism}\}
$$
is equal to
$$
S' = \{f \in \text{Arrows}(\mathcal{C}) \mid
\text{there exist }g, h\text{ such that }gf, fh \in S\}
$$
and is the smallest saturated multiplicative system containing $S$.
In particular, if $S$ is saturated, then $\hat S = S$.
\end{lemma}

\begin{proof}
It is clear that $S \subset S' \subset \hat S$ because elements of
$S'$ map to morphisms in $S^{-1}\mathcal{C}$ which have both left
and right inverses. Note that $S'$ satisfies MS4, and that
$\hat S$ satisfies MS1. Next, we prove that $S' = \hat S$.

\medskip\noindent
Let $f \in \hat S$. Let $s^{-1}g = ht^{-1}$ be the inverse morphism
in $S^{-1}\mathcal{C}$. (We may use both left fractions and right
fractions to describe morphisms in $S^{-1}\mathcal{C}$, see
Lemma \ref{lemma-multiplicative-system}.)
The relation $\text{id}_X = s^{-1}gf$ in $S^{-1}\mathcal{C}$ means
there exists a commutative diagram
$$
\xymatrix{
 & X' \ar[d]^u & \\
X \ar[ru]^{gf} \ar[r]^{f'} \ar[rd]_{\text{id}_X} &
X'' &
X \ar[lu]_s \ar[l]_{s'} \ar[ld]^{\text{id}_X} \\
& X \ar[u]_v &
}
$$
for some morphisms $f', u, v$ and $s' \in S$. Hence $ugf = s' \in S$.
Similarly, using that $\text{id}_Y = fht^{-1}$ one proves that
$fhw \in S$ for some $w$. We conclude that $f \in S'$. Thus
$S' = \hat S$. Provided we prove that $S' = \hat S$ is a
multiplicative system it is now clear that this implies that $S' = \hat S$
is the smallest saturated system containing $S$.

\medskip\noindent
Our remarks above take care of MS1 and MS4, so to finish the proof of the
lemma we have to show that LMS2, RMS2, LMS3, RMS3 hold for $\hat S$.
Let us check that LMS2 holds for $\hat S$. Suppose we have a solid diagram
$$
\xymatrix{
X \ar[d]_t \ar[r]_g & Y \ar@{..>}[d]^s \\
Z \ar@{..>}[r]^f & W
}
$$
with $t \in \hat S$. Pick a morphism $a : Z \to Z'$ such that
$at \in S$. Then we can use LMS2 for $S$ to find a commutative diagram
$$
\xymatrix{
X \ar[d]_t \ar[r]_g & Y \ar[dd]^s \\
Z \ar[d]_a \\
Z' \ar[r]^{f'} & W
}
$$
and setting $f = f' \circ a$ we win. The proof of RMS2 is dual to this.
Finally, suppose given a pair of morphisms $f, g : X \to Y$ and
$t \in \hat S$ with target $X$ such that $ft = gt$.
Then we pick a morphism $b$ such that $tb \in S$. Then
$ftb = gtb$ which implies by LMS3 for $S$ that there exists an $s \in S$
with source $Y$ such that $sf = sg$ as desired. The proof of
RMS3 is dual to this.
\end{proof}








